\section*{Chapter \ref{ch:iv1:statistics} --- Statistics}
\markboth{\protect\omTocMark{}{Answers: Statistics}}{\protect\omTocMark{}{Answers: Statistics}}

\begin{solution}{iq:iv1:statistics:1}
The normal approximation gives $z = 2$, a p-value of about $0.046$, and rejects fairness. The exact binomial p-value is about $0.057$, and it does not reject.

Under a fair coin, heads have mean 50 and standard deviation $\sqrt{100 \cdot \frac14} = 5$\index{normal approximation to the binomial}, which gives $z = 2$. The gap comes from the count being a whole number: the exact tail starts at 60 and includes all of it, so the fair comparison uses $59.5$, which gives $z = 1.9$ and a p-value of $0.057$, in line with the exact one.

\iqwrong{Treating $0.05$ as a sharp truth, and saying that the two methods disagree about the coin. The evidence is borderline in both cases, and that is the honest report.}
\end{solution}

\begin{solution}{iq:iv1:statistics:2}
The interval is $0.1\% \pm 0.098\%$, so about $[0.002\%, 0.198\%]$. Halving the width takes four times the data: 1600 days.

The standard error is $1\%/\sqrt{400} = 0.05\%$\index{standard error of the mean}, and $1.96$ of those is $0.098\%$. Standard errors fall with $\sqrt{n}$, so half the width needs four times the sample.

\iqwrong{Saying that it needs twice the data. The square-root law says four times, and that is why adding precision gets expensive quickly.}
\end{solution}

\begin{solution}{iq:iv1:statistics:3}
About 36\% of the ideas that pass are duds. And a p-value of $0.05$ does not mean a 5\% chance that the null is true.

Out of 1000 ideas\index{false discovery}, 900 are duds, and $0.05$ of them pass: 45. Of the 100 good ideas, 80 pass. So the fraction of duds among those that pass is $45/125 = 0.36$.

A p-value is the chance of data this extreme if the null is true. The chance of the null given the data needs a prior, as this calculation shows. If all ideas were bad, every pass would be a dud, whatever the p-value.

\iqwrong{Saying that a pass gives a 95\% chance that the idea is good. That reverses the conditional, and it ignores how rare good ideas are.}
\end{solution}

\begin{solution}{iq:iv1:statistics:4}
Dividing by $n - 1$ corrects for measuring deviations from the sample mean, rather than the true mean. The plain version with $n$ averages $\frac{n - 1}{n}\sigma^2$, which is $0.8\sigma^2$ for samples of size 5.

Start from
\[
\sum (X_i - \bar{X})^2 = \sum (X_i - \mu)^2 - n(\bar{X} - \mu)^2.
\]
The mean of the first term is $n\sigma^2$, and the mean of the second is $n \cdot \sigma^2/n = \sigma^2$\index{Bessel's correction}. So the sum averages $(n - 1)\sigma^2$. The sample mean sits closer to the data than the true mean does, so deviations from it are too small. This does not need normality, and a simulation on exponentials confirms it.

\iqwrong{Saying that you lose a data point. The term ``degree of freedom'' is a good name, but you should be able to show the algebra.}
\end{solution}

\begin{solution}{iq:iv1:statistics:5}
At least one passes with chance $1 - 0.95^{20} \approx 0.64$. A level of about $0.0025$ per test keeps the chance at about 5\%.

Exactly, the per-test level is $1 - 0.95^{1/20} = 0.00256$\index{multiple testing}. Bonferroni's simple rule, $0.05/20 = 0.0025$, gives a family chance of $0.049$. It holds even when the tests are dependent, because it only uses the union bound.

\iqwrong{Thinking that 5\% stays 5\%, because each test has a 5\% error. The chance that at least one is wrong rises quickly with the number of tests.}
\end{solution}

\begin{solution}{iq:iv1:statistics:6}
About 31 of them, since $1000/2^5 = 31.25$. Screening for a clean five-year record finds those 31, and it tells you nothing about skill.

Each manager has a chance of $1/32$ of a clean record\index{survivorship bias}. So a clean record is not rare in a large group, and the investor would see a list of 31 winners without any skill existing at all. Next year, each still has a 50\% chance.

If you see a clean record, the right question is how many people started. That number is usually not on the page, because the losers have closed.

\iqwrong{Saying that $1/32$ is about 3\%, so the record is significant at the 5\% level. That test applies to one manager chosen in advance, not to the best of a thousand.}
\end{solution}

\begin{solution}{iq:iv1:statistics:7}
The exact two-sided p-value is $0.0414$, using a sign test.

If the median is zero, each day is up with a chance of $\tfrac12$, whatever the distribution\index{sign test}. So the count of up days is binomial with $n = 20$ and $p = \tfrac12$. Then
\[
P(X \ge 15) = \frac{\binom{20}{15} + \dots + \binom{20}{20}}{2^{20}} = \frac{21\,700}{1\,048\,576},
\]
and doubling this gives the two-sided p-value. A simulation with Cauchy draws gives the same chance, as it should. The test ignores sizes, so it loses power on normal data, but it does not break on fat tails.

\iqwrong{Using a $t$ test on the returns. That needs the sizes, which you were not given, and it would be unreliable on 20 fat-tailed points.}
\end{solution}

\begin{solution}{iq:iv1:statistics:8}
The p-value is $1/20 = 0.05$. That is the smallest that any permutation test can give with two groups of 3.

If the labels have no effect, each choice of which 3 of the 6 values are called A is equally likely\index{permutation test}. There are $\binom{6}{3} = 20$ choices. The observed one, $\{5, 7, 9\}$, has the largest possible sum, 21, and it is the only one to reach it. So the p-value is $1/20$.

This is the floor for two groups of 3, since a single assignment is the most extreme you can have. At this sample size, you can never get below 5\%, whatever the actual size of the effect.

\iqwrong{Running a $t$ test, which gives a much smaller value. It gets that from an assumption that the values are normal, which you cannot check with 6 points.}
\end{solution}

\begin{solution}{iq:iv1:statistics:9}
The normal interval is $1 \pm 0.877$, and the $t$ interval is $1 \pm 1.241$. The difference comes from the standard deviation being estimated, and it stops mattering at about 30 points, where the $t$ value is $2.05$.

The standard error is $1/\sqrt{5} = 0.447$\index{Student's t}. Multiplying by $1.96$ and by $2.776$ gives the two intervals. When $\sigma$ is estimated from 5 points, a low estimate is common, which makes the ratio large more often than a normal variable would be. So tails are fatter. A simulation shows the normal interval covering only about $87.8\%$ of the time, and the $t$ interval $95\%$.

\iqwrong{Saying that the difference is about non-normal data. Here the data are normal. The extra width comes only from not knowing $\sigma$.}
\end{solution}

\begin{solution}{iq:iv1:statistics:10}
The rule of thumb is $1/\sqrt{n} \approx 0.099$ near zero. The 95\% interval for $r = 0.5$ is about $[0.34, 0.63]$.

Fisher's transformation\index{Fisher's z transformation}, $z = \operatorname{arctanh}(r)$, makes the sampling law close to normal, with a standard error of $1/\sqrt{n - 3} = 0.1$, whatever the true correlation. So $z = 0.549 \pm 0.196$. Transforming back with $\tanh$ gives the interval. A simulation with a true correlation of $0.5$ confirms a standard deviation of $0.1$ on the transformed scale, and 95\% coverage.

The interval is not symmetric around $0.5$, because correlation is bounded by 1, and the transformation stretches values near $\pm 1$. An estimate of $0.5$ has more room below it than above.

\iqwrong{Using $0.5 \pm 2 \cdot 0.099$. Near zero it is fine, but at $0.5$ the true standard error is about $\frac{1 - r^2}{\sqrt{n}} \approx 0.074$, which is much smaller, and the interval is skewed.}
\end{solution}

\begin{solution}{iq:iv1:statistics:11}
The exact bound is $1 - 0.05^{1/100} \approx 0.0295$, and the rule of thumb is $3/n = 0.03$. The naive estimate is 0, which claims that failure is impossible.

Find the largest failure chance $p$ that is still consistent with no failures\index{rule of three}. You need $(1 - p)^{100} \ge 0.05$. Taking logs gives $100p \approx -\ln 0.05 \approx 3$. So 100 clean days say that the chance is below about 3\% a day. That is a chance of a failure about every 33 days, which is hardly reassuring.

\iqwrong{Giving the point estimate 0, or $1/100$. An interval is needed, and with no failures the upper end is what matters.}
\end{solution}

\begin{solution}{iq:iv1:statistics:12}
You need 16 years at a Sharpe ratio of $0.5$, 4 years at 1 and 1 year at 2.

Measured over $T$ years, an annual Sharpe ratio has a standard error near $1/\sqrt{T}$\index{Sharpe ratio inference}. This is the same as the standard error of a mean in units of its own standard deviation, and it does not depend on how often you sample. So the ratio is $\mathrm{SR}\sqrt{T}$, and it reaches 2 when $T = 4/\mathrm{SR}^2$. A simulation with 16 years of daily returns confirms a standard error of $0.25$.

So a three-year record cannot separate a Sharpe of $0.5$ from zero, and no realistic record can do so for a slow strategy. Only high-frequency strategies with very high Sharpe ratios show their skill quickly.

\iqwrong{Thinking that daily data help, because there are so many points. A Sharpe ratio is a mean over a volatility. Only the volatility benefits from the finer sampling.}
\end{solution}

\begin{solution}{iq:iv1:statistics:13}
You need $n = (1.96 + 0.84)^2(\sigma/\delta)^2 = 3136$ days, about $12.4$ years. A doubled effect needs only a quarter: 784 days, about $3.1$ years.

The test rejects when the sample mean is more than $1.96$ standard errors from zero\index{power calculation}. To have 80\% power, the true effect must be $1.96 + 0.84$ standard errors from zero, because $0.84$ is the 80\% point of a normal. So $\delta\sqrt{n}/\sigma = 2.8$, and squaring gives $n = 7.84 \cdot 400$. A simulation of the test confirms 80\% power.

\iqwrong{Using only the $1.96$, which gives about 1537 days. At that size you have just a 50\% chance of detecting the effect, because half of the time the estimate falls below the true value.}
\iqwhy{Plan your sample before you test. A daily edge of 5 basis points in 1\% volatility is a Sharpe near 0.8, and it takes more than a decade to confirm by this rule.}
\end{solution}

\begin{solution}{iq:iv1:statistics:14}
Bonferroni rejects one test, and Benjamini--Hochberg rejects three. Bonferroni controls the chance of any false rejection, and Benjamini--Hochberg controls the expected fraction of false rejections among the rejections.

Bonferroni compares each p-value with $0.05/10 = 0.005$. Only $0.001$ passes.

Benjamini--Hochberg\index{Benjamini--Hochberg procedure} sorts them, and finds the largest $k$ with $p_{(k)} \le k \cdot 0.05/10$. Here $0.001 \le 0.005$, $0.008 \le 0.010$, $0.012 \le 0.015$, but $0.041 > 0.020$, and no later one passes. So the first three are rejected. If you are searching for signals, and will validate them after, a false discovery rate is the natural choice. If one false result is costly, use the family-wise rule.

\iqwrong{Stopping at the first failure in a step-up rule without looking further. Benjamini--Hochberg takes the last $k$ that passes, even if an earlier one fails.}
\end{solution}

\begin{solution}{iq:iv1:statistics:15}
The method of moments gives $\hat\theta = 8/4 = 2$ and $\hat{k} = 4/2 = 2$. Maximum likelihood is more efficient, especially for the shape, at the cost of a numerical solution.

Match the moments\index{method of moments}. Variance over mean gives $\theta$, and then $k$ is mean over $\theta$.

The method is simple, and it gives a good starting value. But it uses the sample variance, which is noisy for skewed data. A fit with 300 samples of size 50 shows a lower variance for the maximum-likelihood shape than for the moment shape. Maximum likelihood is asymptotically efficient\index{maximum likelihood}, if the model is right. So it is common to start from the moment values and then iterate.

\iqwrong{Saying that the two methods are the same, since they use the same data. They use it differently, and they agree only for special families such as the Poisson or the normal.}
\end{solution}

\begin{solution}{iq:iv1:statistics:16}
The maximum-likelihood estimate is $\hat\lambda = 1/\bar{X}$. It is biased high: $E[\hat\lambda] = \frac{n}{n - 1}\lambda$. The unbiased version is $\frac{n - 1}{n\bar{X}}$, which is $0.4$ against $0.5$ for $n = 5$ and $\bar{X} = 2$.

The log-likelihood is $n\ln\lambda - \lambda\sum x_i$\index{maximum likelihood}. Setting its derivative to zero gives $\lambda = n/\sum x_i$.

The sum is Gamma with shape $n$, so $E[1/\sum x_i] = \lambda/(n - 1)$, which gives the bias. In words: $1/x$ is convex, so by Jensen's inequality the average of $1/\bar{X}$ is above $1/E[\bar{X}]$\index{Jensen's inequality}. A simulation with rate 3 and $n = 5$ gives $3.75$ for the first and 3 for the second.

\iqwrong{Thinking that an unbiased mean gives an unbiased rate. A nonlinear function of an unbiased estimator is biased.}
\end{solution}

\begin{solution}{iq:iv1:statistics:17}
The Cram\'er--Rao bound is $\lambda/n$, and $\mathrm{Var}(\bar{X}) = \lambda/n$, so the bound is reached. At $\lambda = 2$, the asymptotic variances for the chance of zero are $0.0366/n$ for $e^{-\bar{X}}$ and $0.117/n$ for the fraction of zero days. So $e^{-\bar{X}}$ is about $3.2$ times better, if the data really are Poisson.

The Fisher information per day is $1/\lambda$\index{Cram\'er--Rao bound}, since the score is $x/\lambda - 1$ and its variance is $\lambda/\lambda^2$. So the bound is $\lambda/n$.

For $e^{-\bar{X}}$, use the delta method: the slope is $-e^{-\lambda}$, which gives $\lambda e^{-2\lambda}/n$. The fraction of zeros is a binomial proportion, with variance $e^{-\lambda}(1 - e^{-\lambda})/n$. Simulations with 200 days confirm both.

The first needs the model. With overdispersed counts of mean 2, and a true chance of a zero of $1/3$, it would report $e^{-2} \approx 0.135$, far off, while the fraction stays correct. Efficiency has a price, which is the model assumption.

\iqwrong{Preferring the fraction because it is ``non-parametric'', and calling that better without qualification, or the reverse. Each one is the right choice under a different view of how much you trust the model.}
\end{solution}

\begin{solution}{iq:iv1:statistics:18}
About $1 - 1/e \approx 63.2\%$. The bootstrap estimates the sampling variability of a statistic, by treating the sample as if it were the population.

A given point is missed by one draw with chance $1 - 1/n$\index{bootstrap}, and by all $n$ with chance $(1 - 1/n)^n \to 1/e$. So $63.2\%$ of the points appear, some of them repeatedly.

The idea: you would like to draw many samples from the true distribution, and see how the statistic varies. You cannot do that, so you draw from the sample itself, which is the best estimate you have of that distribution. For a mean, the bootstrap standard error matches the formula $s/\sqrt{n}$, as a simulation confirms. For a median or a Sharpe ratio, where there may be no simple formula, it gives you one with no algebra.

\iqwrong{Saying that the bootstrap creates more data. It creates no information about the population. It only shows how much the statistic moves under resampling.}
\end{solution}

\enlargethispage{-1\baselineskip}%
\begin{solution}{iq:iv1:statistics:19}
The variance is 3 times larger: $(1 + \rho)/(1 - \rho) = 3$, so the standard error is $\sqrt{3} \approx 1.73$ times larger. A plain bootstrap reports the naive figure. Use a block bootstrap.

The variance of a mean is\index{autocorrelation and the variance of the mean}
\[
\frac{\sigma^2}{n}\Bigl(1 + 2\sum_{k \ge 1} \rho^k\Bigr) = \frac{\sigma^2}{n} \cdot \frac{1 + \rho}{1 - \rho}.
\]
With $\rho = 0.5$, the sum is $1 + 2 \cdot 1 = 3$. Paths of length 2000, simulated, show this factor.

A plain bootstrap draws points independently, so it destroys the autocorrelation and reports $\sigma^2/n$, as a simulation on one path confirms. A block bootstrap\index{block bootstrap} draws runs of neighbouring points, longer than the memory of the series, here blocks of 50, and that recovers most of the factor.

\iqwrong{Using the plain bootstrap as an all-purpose fix. It assumes independent draws, and the result is too narrow on any autocorrelated data.}
\iqfollow{And for negative autocorrelation?}{The factor is below 1. At $\rho = -0.5$, it is $1/3$, so the mean is known better than the naive formula says.}
\end{solution}

\begin{solution}{iq:iv1:statistics:20}
Under normality, the standard errors are $\sqrt{6/n} \approx 0.077$ for skewness and $\sqrt{24/n} \approx 0.155$ for excess kurtosis. Jarque--Bera is $\frac{1000}{6}\bigl(0.25 + \frac94\bigr) \approx 417$, compared with a chi-square with 2 degrees of freedom, whose 95\% point is $5.99$.

The statistic is the sum of the two squared $z$-scores\index{Jarque--Bera test}, $\frac{n}{6}S^2 + \frac{n}{24}K^2$. Simulations of normal samples of 1000 confirm both standard errors.

The test rejects easily with financial data, and that is its weakness. With many points it finds any departure, however small. And it says only that the data are not normal, not why. Stochastic volatility with normal shocks fails it too, in the simulation more than 99\% of the time, so a rejection does not imply fat-tailed shocks. Skewness and kurtosis estimates are also very sensitive to a few extreme days, so you should look at those before you trust them.

\iqwrong{Treating rejection as a useful finding. For daily returns the test almost always rejects, so a chart of the tails tells you more.}
\end{solution}

\begin{solution}{iq:iv1:statistics:21}
If normal, the chance is about $2.87 \times 10^{-7}$ a day, about once in $13\,800$ years. If Student $t$ with 3 degrees of freedom, the chance is about $0.0016$ a day, about once in $2.5$ years. Such a day is not a miracle. It tells you that the model is wrong.

For the normal, $\Phi(-5) \approx 2.87 \times 10^{-7}$, and with 252 trading days per year you get the first number.

For the $t$, the variance is $\nu/(\nu - 2) = 3$\index{fat tails}, so a rescaled move of $-5$ is a raw $t$ value of $-5\sqrt{3} = -8.66$. Its tail is about $0.0016$. A simulation confirms that tail.

You should not conclude that a very rare thing has happened. Compare the likelihood under each model. If 5-sigma days occur every few years, then the normal is the model to drop, and a stress limit built on it has to be redone.

\iqwrong{Saying that the event is impossible, or that it is simply bad luck. A probability as low as that is a statement about a model, and evidence that it is wrong is much more likely than the event itself.}
\end{solution}

\begin{solution}{iq:iv1:statistics:22}
The chi-square statistic is 8, below $11.07$, so fairness is not rejected (p about $0.156$). You could test face 1 alone, with $z \approx 2.45$ (p about $0.014$), but you chose it after looking. So you must correct for the 6 faces, which gives about $0.09$.

Each expected count is 20\index{chi-square test}. The statistic is
\[
\sum \frac{(O - 20)^2}{20} = \frac{100 + 4 + 0 + 4 + 16 + 36}{20} = 8.
\]
For one face, the count is binomial, with standard deviation $\sqrt{120 \cdot \frac16 \cdot \frac56} = 4.08$, so $z = 10/4.08 = 2.45$.

\iqwrong{Testing the largest face and reporting $0.014$. It is the largest of six by construction, so that p-value is six times too good.}
\end{solution}

\begin{solution}{iq:iv1:statistics:23}
The mean is $50.5$ and the standard deviation is about $4.97$. Thirty runs is $4.1$ standard deviations low, which suggests momentum or clustering. Sixty-five is $2.9$ high, which suggests reversal.

The number of runs is 1 plus the number of changes\index{runs test}. Each of the 99 neighbouring pairs is a change with a chance of $\tfrac12$, and these are independent for a fair sequence. So the count of changes is binomial with $n = 99$, with mean $49.5$ and variance $99/4$. A simulation confirms both.

Note that the test looks at order, not at the proportion of up days. A market can have a positive average return and still look perfectly random on this test.

\iqwrong{Thinking that a long run is suspicious by itself. Long runs are common in fair data, and people tend to underestimate them. What counts is the number of runs overall, against this mean and standard deviation.}
\end{solution}

\begin{solution}{iq:iv1:statistics:24}
As independent samples $z \approx 0.71$, and as pairs $z \approx 2.24$. Pairing is right, because both algorithms traded the same orders.

As independent samples, the standard error of the difference is $\sqrt{(100 + 100)/100} = 1.41$. As pairs, the difference on each order has variance $2 \cdot 100 \cdot (1 - 0.9) = 20$\index{paired test}, so the standard error is $\sqrt{20/100} = 0.447$. That is $\sqrt{10} \approx 3.16$ times smaller. A simulation of correlated pairs gives the same figure.

The orders' difficulty is shared, so it cancels in the difference. To compare methods, design the test so that shared noise can cancel: run both on the same orders, or on randomly assigned neighbouring orders. Here the same data move from ``no evidence'' to ``significant'' purely by using the right design.

\iqwrong{Using the independent test because that is the default, and concluding that the two algorithms cannot be told apart. The formula has to reflect how the data were made.}
\end{solution}

\begin{solution}{iq:iv1:statistics:25}
Yes, it is significant. The difference is 3 with a standard error of $\sqrt{2}$, so $z \approx 2.12$ and p is about $0.034$. With equal errors, intervals overlap until the gap reaches $3.92$ standard errors, but significance needs only $2.77$.

The two intervals are $0 \pm 1.96$ and $3 \pm 1.96$, which overlap between $1.04$ and $1.96$\index{confidence interval overlap}. But the standard error of a difference is not the sum of the two standard errors. It is $\sqrt{1 + 1}$. So non-overlap is a very strict test, with a real level of about $0.006$. Overlap does not mean ``no difference''.

In general, with standard errors $s_1$ and $s_2$, the intervals overlap while the gap is below $1.96(s_1 + s_2)$, and the difference is significant once the gap exceeds $1.96\sqrt{s_1^2 + s_2^2}$. So a significant difference with overlapping intervals is always possible, and the two thresholds differ by a factor of up to $\sqrt{2}$, reached when $s_1 = s_2$.

\iqwrong{Saying that overlapping intervals mean no significant difference. Non-overlap is the stricter rule, by a factor of up to $\sqrt{2}$ in the gap.}
\end{solution}

\begin{solution}{iq:iv1:statistics:26}
Pearson is about $0.59$ and Spearman is 1. Spearman gives the fairer summary of this relation, which is perfectly increasing.

Spearman is the correlation of the ranks\index{rank correlation}. The ranks of $y$ are those of $x$, so it is exactly 1. Pearson measures a linear relation, so one point far off the line pulls it down. Without that point it is 1 as well.

Use Spearman when you care about ordering, or when you have outliers or fat tails. It ignores sizes, so it may miss the fact that the top point really matters. If you size a position by a signal, the actual linear relation matters, and Pearson on winsorised data may be more relevant. Show both, and if they differ look for the reason.

\iqwrong{Saying that $0.59$ is the true strength of the relation, or dropping the point because it looks odd. Look at why it is there first.}
\end{solution}

\begin{solution}{iq:iv1:statistics:27}
The best divisor is $c = n + 1$. For $n = 10$, the mean squared errors are $0.2222$ for $c = 9$, $0.19$ for $c = 10$ and $0.1818$ for $c = 11$.

Let $S = \sum(X_i - \bar{X})^2$, so $S/\sigma^2$ is chi-square with $m = n - 1$ degrees of freedom\index{mean squared error}. So $E[S] = m\sigma^2$ and $E[S^2] = (2m + m^2)\sigma^4$. Then, in units of $\sigma^4$,
\[
E\Bigl[\Bigl(\frac{S}{c} - 1\Bigr)^2\Bigr] = \frac{2m + m^2}{c^2} - \frac{2m}{c} + 1.
\]
Minimising over $1/c$ gives $c = \frac{2m + m^2}{m} = m + 2 = n + 1$. The minimum value is $2/(n + 1)$.

The best divisor gives a biased estimate: it averages $\frac{9}{11}\sigma^2$ here. That is a good trade, since variance estimates have high variance, and shrinking them reduces that\index{bias--variance trade-off}. Unbiasedness is one property, and it is not the same as accuracy. The result depends on the normal kurtosis. With a Laplace, a simulation shows a best divisor above 12, so shrink even more on fat tails.

\iqwrong{Saying that $n - 1$ must be best, as it is unbiased. Unbiased means right on average. You can have a smaller mean squared error by accepting a small bias, even though the unbiased estimate is the closer one on most samples: here $S/11$ beats $S/9$ on only about 36\% of them.}
\end{solution}

\begin{solution}{iq:iv1:statistics:28}
Reject when $\bar{X} > \tfrac12$, which is the likelihood-ratio test. You need $n = 11$, as $(2 \cdot 1.645)^2 = 10.8$.

By the Neyman--Pearson lemma\index{Neyman--Pearson lemma}, the most powerful test at any level uses the likelihood ratio. For two normals with the same variance,
\[
\ln\frac{L_1}{L_0} = \sum_i \Bigl(x_i - \frac12\Bigr) = n\Bigl(\bar{x} - \frac12\Bigr),
\]
so the test depends only on the sample mean. Setting the cut-off at $\tfrac12$ gives equal error rates, by symmetry.

Each error is $P\bigl(Z > \tfrac12\sqrt{n}\bigr)$, because the cut-off is half a mean-gap away, in units of $1/\sqrt{n}$. For this to be at most 5\%, you need $\tfrac12\sqrt{n} \ge 1.645$. At $n = 11$ the error is $0.049$, and at $n = 10$ it is above 5\%. A simulation at $n = 11$ confirms this for both.

For a gap of $\delta$ standard deviations, $n \approx 10.8/\delta^2$. Halving the gap takes four times the data.

\iqwrong{Using a one-sided test against the first hypothesis at 5\%, with a cut-off of $1.645/\sqrt{n}$. That fixes only one error rate, and it says nothing about the other. A fair decision between two exact hypotheses controls both.}
\end{solution}

\begin{solution}{iq:iv1:statistics:29}
The real false positive rate is about 19\%. With no limit on looks, it creeps up towards 100\%. A cut-off of about $2.56$ per look restores 5\%, and this is the idea behind group sequential testing, with Pocock boundaries.

Each look is correlated with the previous ones, so the risk does not add up fully, but it does accumulate\index{optional stopping in experiments}. Simulating the $z$-scores of a cumulative sum over 10 looks gives $19\%$. With 5 looks it is $14\%$.

With no limit, the law of the iterated logarithm says that a random walk's $z$-score eventually exceeds any fixed level\index{law of the iterated logarithm}. So you are certain to ``find'' an effect if you are willing to wait long enough. A simulation with 2000 looks is already above 40\%.

A cut-off of $2.56$ per look, which is a nominal p-value of about $0.01$, keeps the total at 5\% for 10 looks. Other schemes spend the error unevenly, so that early looks need very strong evidence.

\iqwrong{Thinking that peeking is harmless if you do not change the procedure. The stopping decision is part of the procedure, and a p-value computed as though the sample size were fixed is wrong.}
\iqwhy{In research and in backtesting alike, stopping when results look good is the same trap as trying many signals. Decide the sample before you start, or use a boundary designed for stopping.}
\end{solution}

\begin{solution}{iq:iv1:statistics:30}
The bootstrap maximum equals the sample maximum with probability $1 - (1 - 1/n)^n \to 1 - 1/e \approx 0.632$. The true sampling law has no such lump. Use the exact law: $P(\max/\theta \le u) = u^n$.

A bootstrap sample contains the sample's largest point unless all $n$ draws miss it\index{bootstrap failure}, so the bootstrap distribution has an atom of $63.2\%$ at the sample maximum. Any bootstrap interval is therefore degenerate on one side. It cannot reflect that the true $\theta$ lies above every observed point.

The true law is continuous. $n(\theta - \max)/\theta$ tends to an exponential with mean 1, and simulations confirm its mean and a tail value. Since $\max/\theta$ has the exact law above, a 95\% one-sided interval is $\theta \le \max/0.05^{1/n}$, and for $n = 1000$ the factor is $1/0.9970$.

Why it fails: the bootstrap works when the statistic is a smooth function of the empirical distribution. An extreme value depends on what happens at the edge, and the sample holds no information from beyond the edge. A subsample of size $m \ll n$ works better, as does a parametric bootstrap.

\iqwrong{Saying that the bootstrap always works if you take enough resamples. More resamples only reduce the Monte Carlo error. They do not repair a method that is wrong for this statistic.}
\iqwhy{Tail risk measures are close to this case. Bootstrapping the worst day or the extreme quantile of a return series understates how bad things can be, for the same reason.}
\end{solution}

\begin{solution}{iq:iv1:statistics:31}
Under normality the standard error is about $0.093$. With a skewness of $-1$ and a kurtosis of 6, it is about $0.109$, so 16\% wider.

Write $\hat{s} = \bar{X}/\hat\sigma$. Under normality, $\bar{X}$ and $\hat\sigma$ are independent, with variances $\sigma^2/n$ and about $\sigma^2/2n$. The delta method gives\index{delta method}
\[
\mathrm{Var}(\hat{s}) \approx \frac{1}{n} + s^2 \cdot \frac{1}{2n} = \frac{1 + s^2/2}{n}.
\]
With $s = 0.3$ and $n = 120$, that gives $0.0933$.

With the given formula, the numerator is $1 + 0.3 + \frac54 \cdot 0.09 = 1.41$, so the result is $0.1085$. The formula is also tested by a simulation with a mixture of skewed, fat-tailed returns, where it is within 5\%.

Negative skew makes it wider, because it couples the mean to the volatility: a large loss both lowers the mean and raises the volatility, so errors in the two move the ratio in the same direction. Fat tails make the volatility estimate noisier, so the second term is larger.

\iqwrong{Applying $1/\sqrt{n}$ with no correction. For a low Sharpe and thin tails this is fine. On a strategy that sells options, with heavy negative skew, the interval is much too narrow.}
\iqfollow{What about annual units?}{Multiply by $\sqrt{12}$: an annual Sharpe of $1.04$ has a standard error of about $0.32$ for normal returns and about $0.38$ in the fat-tailed case. The record is then about 3 standard errors from zero.}
\end{solution}

\begin{solution}{iq:iv1:statistics:32}
The correct annual Sharpe ratio is about $0.86$, not $1.04$. With autocorrelation of $-0.2$, it is about $1.25$. Illiquid or marked-to-model strategies often show positive autocorrelation, because their reported prices move slowly.

The annual return is the sum of 12 monthly returns, so its mean is $12\mu$\index{Sharpe ratio annualisation}, and its variance is
\[
\sigma^2\Bigl(12 + 2\sum_{k = 1}^{11}(12 - k)\rho^k\Bigr).
\]
The usual rule keeps only the first term. With $\rho = 0.2$ the bracket is about $17.4$, and so the annual Sharpe ratio is $0.3 \cdot 12/\sqrt{17.4} \approx 0.86$. A simulation of 100\,000 years of this process gives the same.

Positive autocorrelation makes a strategy's monthly volatility look lower than its risk over a long horizon, and therefore a Sharpe ratio that is too high. Stale prices cause this: a private asset valued by appraisal, or an illiquid book marked to a model, moves slowly, which smooths returns month to month. Real momentum can also contribute.

\iqwrong{Applying $\sqrt{12}$ in all cases. It assumes independent returns. If you see smooth returns from an illiquid book, you should treat the autocorrelation as a warning and unsmooth the series before computing risk.}
\end{solution}

\begin{solution}{iq:iv1:statistics:33}
The standard error of the difference is about $\sqrt{2(1 - 0.8)/5} \approx 0.283$, against about $0.632$ if independent. So the correlation makes the comparison $2.2$ times sharper, but a gap of $0.28$ is still only one standard error.

Each measured Sharpe ratio has a variance of about $1/T$. The covariance of two of them is about $\rho/T$\index{comparing Sharpe ratios}, since it comes mainly from the correlation between the two means. So the difference has a variance of $2(1 - \rho)/T$. With daily data, the volatility term is negligible, because each day's Sharpe ratio is small, and a simulation with 5 years of daily data gives $0.283$.

So over 5 years, two equally good strategies with this correlation will often appear different: a gap above $0.28$ in favour of one of them occurs more than 10\% of the time, in the simulation. If you have just picked the better of two similar strategies on the basis of a backtest, you probably picked noise. The choice may still be harmless, as they are correlated, but you should not expect the gap to persist.

\iqwrong{Using the two separate standard errors, or the one-sample formula $1/\sqrt{T} = 0.447$ for the gap. Both overstate the error of the difference, and they suggest that the strategies cannot be told apart even when they can.}
\iqwhy{An uplift in a backtest from changing a parameter is a difference between two highly correlated strategies. It often has more precision than you might think, but not enough to justify most small improvements.}
\end{solution}

\begin{solution}{iq:iv1:statistics:34}
The sample quantile has a standard error of about $0.118\sigma$, and $2.326\hat\sigma$ has one of about $0.052\sigma$. The first is about $2.3$ times larger.

Take the count of points below the level $q$\index{standard error of a quantile}: it is binomial, with a standard deviation of $\sqrt{np(1 - p)}$. The local slope of its expected value is $nf(q)$, so converting from count to position gives
\[
\mathrm{se}(\hat{q}) = \frac{1}{f(q)}\sqrt{\frac{p(1 - p)}{n}} = \frac{\sqrt{0.0099/1000}}{0.0267} \approx 0.118.
\]
For the other: $\mathrm{se}(\hat\sigma) = \sigma/\sqrt{2n}$, multiplied by $2.326$. Simulations confirm both.

The first is larger because it uses only points near the 1\% tail. Only about 10 points fall there, whereas $\hat\sigma$ uses all of them. In a density-poor tail, each point moves the quantile a lot.

But the second is right only if the normal is. On $t$ returns with 3 degrees of freedom, also tested, the normal rule is biased by more than $0.3\sigma$, whereas the sample quantile is close to unbiased. So choose the sample quantile when you distrust the shape, and report its wide interval. A precise answer to the wrong question is worse.

\iqwrong{Choosing a rule by its standard error alone. The total error includes bias, and at the tail, model error usually dominates sampling error.}
\end{solution}

\begin{solution}{iq:iv1:statistics:35}
For normal returns, the standard error is about $0.141\%$. With a kurtosis of 6, it is about $0.222\%$, which is $\sqrt{5/2} \approx 1.58$ times larger. A chi-square interval has the wrong width, and it undercovers.

The variance of the sample variance is about $\sigma^4(\kappa - 1)/n$, with $\kappa$ the kurtosis\index{standard error of volatility}. A normal has $\kappa = 3$, so this is $2\sigma^4/n$. For the standard deviation, the delta method divides by $2\sigma$:
\[
\mathrm{se}(s) \approx \frac{\sigma}{2}\sqrt{\frac{\kappa - 1}{n}}.
\]
For $\kappa = 3$, this is $\sigma/\sqrt{2n} = 2\%/\sqrt{202} = 0.141\%$. For $\kappa = 6$, it is $1\% \cdot \sqrt{5/101} = 0.222\%$. Simulations with normal and Laplace draws, which have a kurtosis of 6, agree.

The chi-square interval is built on $\kappa = 3$. With fat tails, the actual coverage of a nominal 95\% interval is below 85\% in the simulation. It needs correcting, with a bootstrap or a formula based on the sample kurtosis, and with the caution that the kurtosis estimate is itself very noisy.

\iqwrong{Using $\sigma/\sqrt{2n}$ regardless. The formula looks universal, because it does not contain any shape parameter, but that is because it has assumed one.}
\end{solution}

\begin{solution}{iq:iv1:statistics:36}
The true correlation is about $0.057$, nearly three times the measured $0.02$.

Suppose that $A = S + \text{noise}$ and $B = F + \text{noise}$, with independent noise\index{attenuation}. The noise does not change the covariance, but it does raise the variances. So
\[
\rho_{AB} = \rho_{SF}\sqrt{R_A R_B},
\]
with $R = \mathrm{Var}(\text{true part})/\mathrm{Var}(\text{total})$. Here $R_A = 0.5$ and $R_B = 0.25$, so $\rho_{SF} = 0.02/\sqrt{0.125} \approx 0.057$. A simulation with 40 million points recovers $0.02$ from the true $0.057$. With no noise, the two are the same, which is a check.

In finance, returns are mostly noise, so a raw correlation of $0.02$ to $0.05$ may reflect a strong relation with the small forecastable part. This also tells you that cleaning the signal is valuable: if you halved its noise, the measured correlation would increase. It does not tell you that you can trade better than $0.02$, as you can only trade the noisy signal against the noisy return.

\iqwrong{Treating this correction as a way to boost your measured performance. You can only trade what you measure. The corrected number tells you about the limit, and about what cleaner data could give.}
\end{solution}

\begin{solution}{iq:iv1:statistics:37}
The annual mean has a standard error of about $6.3$ points, at any sampling frequency. The annual volatility has one of about $0.28$ points with daily data, $1.29$ with monthly data, and it would keep falling with finer sampling. So you can estimate volatility far better than the mean, roughly 22 times better here. In practice, microstructure noise limits how fine you can go.

For the mean: the sum of returns over $T$ years depends only on the start and end prices\index{mean versus variance estimation}, so sampling more often adds no information. Its standard error is $\sigma/\sqrt{T} = 20/\sqrt{10}$.

For the volatility: with $n$ returns, the standard error is $\sigma/\sqrt{2n}$. For $n = 252 \cdot 10 = 2520$, that is $20/\sqrt{5040} = 0.28$. Simulations with 120 and 2520 steps confirm both.

So estimating means needs long histories. Estimating volatility needs fine sampling.

At very fine scales, the price is not a pure diffusion. Bid-ask bounce and discrete ticks add noise, which grows relative to the variance, and returns have jumps and intraday patterns too. So there is a sampling interval below which the estimate gets worse. Also, volatility itself changes, so a long sample does not answer the question of what it is today.

\iqwrong{Saying that more data gives a better estimate of everything. Time span and sampling frequency are different kinds of ``more'', and they help different quantities.}
\end{solution}

\begin{solution}{iq:iv1:statistics:38}
The risk is $d - (d - 2)^2 E[1/|X|^2]$, which is below $d$ for every $\theta$. At $\theta = 0$ and $d = 10$ the risk is 2, against 10 for $X$.

Stein's lemma says $E[(X_i - \theta_i)g(X)] = E[\partial_i g(X)]$\index{Stein's lemma}. Write the estimator as $X - h$, with $h = (d - 2)X/|X|^2$. Then
\[
E|X - h - \theta|^2 = d - 2\sum_i E[(X_i - \theta_i)h_i] + E|h|^2 = d - 2E[\nabla \cdot h] + E|h|^2.
\]
A direct calculation gives $\nabla \cdot h = (d - 2)^2/|X|^2$ and $|h|^2 = (d - 2)^2/|X|^2$. So the risk is $d - 2(d - 2)^2 E[1/|X|^2] + (d - 2)^2 E[1/|X|^2]$, which is the formula above. Since $1/|X|^2 > 0$, the risk is below $d$ whatever $\theta$ is, for $d \ge 3$. For $d = 2$, the factor vanishes.

At $\theta = 0$, $|X|^2$ is chi-square, with $E[1/|X|^2] = 1/(d - 2)$, so the risk is $d - (d - 2) = 2$. Simulations at various $\theta$ confirm the formula and the gain\index{James--Stein estimator}.

It is strange, because the components can be anything, unrelated to one another. Yet the total error falls if you shrink them all towards a common point, even if the estimate for each single one is not better. Mean returns on many assets are a good case, because they are noisy and close together. Shrink towards their grand mean instead of towards 0, which leaves a factor with $d - 3$. Portfolio optimisers do better on such estimates.

\iqwrong{Saying this contradicts the sample mean being admissible in one dimension. It holds for $d \le 2$. The result is about the total loss, not any single component.}
\end{solution}

\begin{solution}{iq:iv1:statistics:39}
The Fisher information per observation is $\tfrac12$. So the maximum-likelihood estimate has asymptotic variance $2/n$, the median has $\pi^2/(4n) \approx 2.47/n$, and the sample mean has no finite variance and is no better than one draw, whatever the value of $n$.

The score is $\frac{2z}{1 + z^2}$ with $z = x - \theta$\index{Fisher information}. So
\[
I = \int \frac{4z^2}{(1 + z^2)^2} \cdot \frac{dz}{\pi(1 + z^2)}.
\]
Using $z = \tan\phi$, this becomes $\frac{4}{\pi}\int_{-\pi/2}^{\pi/2} \sin^2\phi\cos^2\phi\, d\phi = \frac{4}{\pi} \cdot \frac{\pi}{8} = \tfrac12$.

The median's variance is $1/(4nf(0)^2)$, with $f(0) = 1/\pi$. So the median has an efficiency of $8/\pi^2 \approx 0.81$ compared with the maximum-likelihood estimate.

The mean of $n$ Cauchy draws has the same Cauchy law as one draw, so averaging gains nothing. A simulation shows that the median of $|\bar{X}|$ stays at 1.

To compute the estimate: the likelihood can have several local peaks, so do not start at random. Start at the median, which is already consistent, and take a few steps of Fisher scoring, adding $\frac{2}{n}\sum \frac{2z_i}{1 + z_i^2}$ at each step\index{Fisher scoring}. That converges to the right peak for large $n$. In a simulation with $n = 301$, the result has a variance close to $2/n$, and below that of the median.

\iqwrong{Using the mean or a variance-based interval. Neither exists in the usual sense for this distribution. Even one big outlier can pull a mean anywhere, which is why a robust estimate or the likelihood itself is needed.}
\end{solution}

\begin{solution}{iq:iv1:statistics:40}
Let $V$ be the number of false rejections and $R$ the total. For one true null with p-value $p_i$, condition on all the other p-values. Let $R_i$ be the number of rejections if $p_i$ were set to 0. Then null $i$ is rejected exactly when $p_i \le R_i\alpha/m$, and in that case the total number of rejections is exactly $R_i$.

The proof has two steps\index{false discovery rate proof}.

Step one: if $p_i \le R_i\alpha/m$, then moving $p_i$ from 0 up to its true value leaves the sorted list passing at $k = R_i$, because $p_i$ is still below the threshold at that rank, and no larger $k$ can pass, since it would also have passed with $p_i = 0$. So there are $R_i$ rejections, including null $i$. If $p_i > R_i\alpha/m$, then null $i$ is not rejected: rejecting it would need a passing $k$ with $k\alpha/m \ge p_i$, so $k > R_i$, and that is again impossible.

Step two: $R_i$ depends only on the other p-values, which are independent of $p_i$, and $p_i$ is uniform. So
\[
E\Bigl[\frac{\mathbf{1}\{i \text{ rejected}\}}{R} \Bigm| \text{others}\Bigr] = \frac{P(p_i \le R_i\alpha/m)}{R_i} = \frac{R_i\alpha/m}{R_i} = \frac{\alpha}{m}.
\]
The $R_i$ cancels, and that is why the result is so clean. Sum over the $m_0$ true nulls:
\[
\mathrm{FDR} = E\Bigl[\frac{V}{\max(R, 1)}\Bigr] = \sum_{i \text{ null}} E\Bigl[\frac{\mathbf{1}\{i \text{ rejected}\}}{R}\Bigr] = \frac{m_0\alpha}{m}.
\]
It is exactly this, not merely at most, under independence. A simulation with $m = 20$, $m_0 = 15$ and $\alpha = 0.1$ gives $0.075$, as stated. A direct test of the first step, setting one p-value over a grid while the others stay fixed, gives a rejection chance of $R_i\alpha/m$.

\iqwrong{Trying to condition on the number of rejections $R$. It depends on $p_i$, so the condition is tangled. The trick is to use $R_i$, which does not depend on $p_i$.}
\end{solution}
